\section{Simu5G: modellazione della rete cellulare 5G}
\label{env:simu5g}

Simu5G è una libreria per OMNeT++ progettata per la valutazione end-to-end di reti 5G, sviluppata come evoluzione della linea di lavoro già avviata con SimuLTE per le reti LTE/LTE-A \cite{virdis2014simulte,nardini2020simu5g}. Simu5G risulta quindi la scelta migliore per sostituire la modellazione LTE presente nel lavoro di riferimento con una modellazione cellulare 5G.

La scelta di Simu5G è motivata da tre fattori principali:
\begin{itemize}
\item integrazione nativa con OMNeT++, che facilita il riuso di workflow e configurazioni basate su moduli, eventi e file \texttt{ini};
\item orientamento alla valutazione end-to-end, utile quando si vogliono osservare effetti di rete che impattano applicazioni e controllo;
\item continuità concettuale con SimuLTE, che rende più naturale il confronto con soluzioni e architetture derivate da ambienti precedentemente basati su LTE.
\end{itemize}

Ai fini del capitolo, è importante evidenziare che Simu5G non viene introdotto come semplice “sostituto tecnologico” nominale di LTE, ma come componente capace di modificare il comportamento osservato del sistema multi-interfaccia. Infatti, la rete cellulare contribuisce alla disponibilità del canale, ai transitori di connettività e alla qualità del flusso informativo che alimenta la logica di controllo e fallback. In una simulazione multi-tecnologia, tali effetti non sono marginali: interagiscono con le altre interfacce e con la dinamica del platoon, influenzando direttamente le metriche sperimentali.

Dal punto di vista della struttura della toolchain, Simu5G si inserisce come framework OMNeT++ aggiuntivo accanto a Veins/PLEXE, condividendo il medesimo motore di simulazione a eventi discreti. Questo rende possibile costruire scenari in cui:
\begin{itemize}
\item la mobilità è governata da SUMO (tramite Veins/TraCI);
\item la logica di platooning e i controllori risiedono in PLEXE;
\item la componente cellulare 5G è modellata da Simu5G.
\end{itemize}
L’integrazione di questi tre livelli è il presupposto tecnico per la parte sperimentale della tesi, che mira a valutare il comportamento della logica importata in una configurazione aggiornata rispetto al lavoro originario.

In questo capitolo ci si limita a inquadrare il ruolo di Simu5G nell’ambiente di sviluppo. Le scelte implementative specifiche (modifiche ai moduli, configurazione radio, integrazione con il progetto di riferimento) saranno discusse nel capitolo dedicato all’integrazione di Simu5G in \texttt{plexe\_hetnet}.
